\chapter{Limitations and Future Enhancements}

\section{Limitations}
\begin{enumerate}[nosep]
\item \textbf{RadHAR is not a subject- or session-independent benchmark.} It has
  two subjects, no subject IDs, and test recordings made minutes after training
  recordings in the same room (Chapter~\ref{ch:leak}). High RadHAR accuracy
  therefore shows generalisation to new recordings, not to new people or rooms.
\item \textbf{MiliPoint}: the radar provides $xyz$ only, action names are not
  published, and the test set has three subjects, so per-subject variance is
  large and only partly captured.
\item \textbf{Synthetic perturbations}: the robustness study perturbs recorded
  point clouds; it does not simulate the physical radar (multipath geometry,
  interference, different mounting heights or rooms).
\item \textbf{Compute}: all experiments ran on a 2-core CPU. Training was capped
  at 15 epochs with early stopping and a 20/24-frame training stride. DGCNN and
  Point Transformer on MiliPoint were run with one seed because of their cost.
  The voxel baseline uses fewer filters than the original RadHAR network.
\item \textbf{Edge claims}: latency was measured on a desktop-class CPU thread and
  storage is computed for INT8 weights; deployment on a micro-controller was not
  carried out.
\item \textbf{Scope}: no fall detection (no labels), single person only, fixed 2-s
  windows, no clinical validation.
\item \textbf{Explainability}: attention maps are not explanations in a causal
  sense \cite{jain2019attention}; occlusion and gradient saliency give
  complementary but approximate views.
\end{enumerate}

\section{Future enhancements}
\begin{enumerate}[nosep]
\item Collect or adopt a dataset with subject IDs, several rooms and fall events;
  evaluate leave-one-subject-out and cross-room.
\item Port the model to an embedded runtime (e.g.\ ONNX Runtime Mobile or an
  MCU inference library) and measure on the radar's companion processor.
\item Physics-based radar simulation for robustness (multipath, occlusion).
\item Self-supervised pre-training on unlabelled radar data (both datasets
  contain many unlabelled frames), and continuous (streaming) recognition.
\end{enumerate}
